\chapter{Grafische Oberfläche im Browser}\label{ch:gui}

\GUI{} ist eine Weboberfläche für \FG. Sie läuft auf einem Rechenserver (am Institut: Server mit NVIDIA RTX 6000
Ada) und wird im Browser bedient. Die GUI rechnet nicht selbst: Sie baut Konfigurationen im Format von
Kapitel~\ref{ch:konfiguration}, ruft die Python-Schnittstelle auf und stellt die Ergebnisse dar. Alles, was in der GUI
entsteht, ist deshalb eine gewöhnliche JSON-Konfiguration, die Sie herunterladen und mit \cmd{fgsolve} oder Python
weiterverwenden können -- und umgekehrt.

\section{Aufbau}
Links steht die Navigation mit den Seiten
\begin{center}
\small
\begin{tabularx}{\linewidth}{lY}
\toprule
Seite & Zweck \\
\midrule
\gui{Editor} & Konfiguration erstellen und bearbeiten (Formulare und JSON), erste Rechnung \\
\gui{Einzelrechnung} & Beispiel oder hochgeladene Datei rechnen, erster Scanpunkt \\
\gui{Scan} & alle Scanpunkte und Sweep-Werte, Diagramme, Vergleich von Läufen \\
\gui{Nahfeld} & Feldschnitte rechnen und darstellen \\
\gui{CAD} & STEP-Dateien hochladen, prüfen, Schnitte ansehen \\
\gui{Optimierung} & Parameter mit exakten Gradienten optimieren \\
\gui{Rechnungen} & alle Rechnungen: Status, Warteschlange, Download, Löschen \\
\gui{System} & Version, Rechenkern, GPU-Zustand, Betrieb \\
\bottomrule
\end{tabularx}
\end{center}
Unter der Navigation stehen Version und Kern von \FG{} (\texttt{cpu} oder \texttt{cuda12}) und die Zahl der laufenden
Rechnungen.

\section{Sitzung und Adresse}
Die Adresszeile enthält Ihre \textbf{Sitzungskennung} \texttt{?s=\ldots}. Speichern Sie die Adresse als Lesezeichen:
Mit ihr kommen Sie -- auch von einem anderen Rechner oder nach einem Neustart des Servers -- zu Ihren Rechnungen und
zur zuletzt bearbeiteten Konfiguration zurück. Wer die Adresse kennt, sieht Ihre Rechnungen; geben Sie sie nur weiter,
wenn das gewollt ist. Ohne Kennung beginnt eine neue Sitzung. Unter \gui{Rechnungen} $\to$ \gui{Meine Sitzung löschen}
entfernen Sie alle Ihre Rechnungen und den Entwurf (nicht rückgängig zu machen).

Ist der Server mit Anmeldung (Proxy) betrieben, gehören Rechnungen zusätzlich Ihrem Benutzernamen; „Nur meine
Rechnungen“ gilt dann über Sitzungen und Geräte hinweg.

\section{Rechnungen im Hintergrund}
Jede Rechnung läuft als \textbf{Hintergrundjob}. Die Oberfläche bleibt bedienbar; der Status aktualisiert sich von
selbst. Standardmäßig rechnet der Server eine Rechnung gleichzeitig (eine GPU, alle Kerne je Rechnung); weitere warten
in einer Warteschlange mit Platznummer. Vor dem Rechnen zeigen Editor und Scan eine \textbf{Abschätzung}: Zahl der
Rechnungen, Matrixgröße $n=4n_f$ und Speicher (Faustwert, Warnung ab 70\,\% des Arbeitsspeichers; bei RK, vielen
Scheiben oder GPU kann der tatsächliche Bedarf weit abweichen).

%-------------------------------------------------------------------------------------------
\section{Editor}
Startseite. Oben wählen Sie die \gui{Quelle}:
\begin{itemize}
  \item \gui{Neu}: ein 1D-Lamellengitter als Vorlage (\gui{Neue Konfiguration}),
  \item \gui{Beispiel}: die Referenzfälle aus \file{examples/} von \FG,
  \item \gui{Datei hochladen}: eine eigene \file{.json}-Konfiguration.
\end{itemize}
Die Konfiguration wird in Reitern bearbeitet:
\begin{description}
  \item[Grundeinstellungen] Dimension, Löser, Fourierordnungen ($N_x$, $N_z$), Perioden, Superstrat und Substrat.
  \item[Einfall] $\lambda$, $\varphi$, $\psi$, $\delta$ jeweils als Zahl, Liste (\texttt{0.5, 0.6, 0.7}; ein einzelnes
        Element als \texttt{0.5,}) oder Bereich \texttt{von:bis:Schritte} (z.\,B.\ \texttt{0.5:0.8:31}).
  \item[Materialien] Art wählen (Brechzahl $n$, Permittivität $\varepsilon$, Hauptwerte mit Euler-Winkeln, Tensor als
        JSON), Material hinzufügen (Name eingeben) und entfernen.
  \item[Schichten] von unten nach oben; mit \gui{$\uparrow$} \gui{$\downarrow$} \gui{$\times$} ordnen bzw.\ löschen. Formulare gibt es für
        \val{homogeneous}, \val{lamellar1d} (Tabelle der Lamellen von/bis/Material), \val{profile1d}, \val{polygon2d},
        \val{contour2d} (Tabelle der Formen: Material, Form, Mitte, Größe, Winkel, Eckpunkte) und \val{profile2d}.
        Andere Typen (CAD, NURBS-Regionen) bearbeiten Sie als JSON der Schicht.
  \item[Löser] RK- und RCWA-Optionen, Rechenkern (Backend), Genauigkeit, Threads.
  \item[JSON] die ganze Konfiguration. \gui{Übernehmen} ersetzt die Formularwerte; \gui{Herunterladen} speichert die
        Datei.
\end{description}
Das Formular ändert nur, was Sie ändern; unbekannte Schlüssel bleiben stehen. \gui{Rechnen} startet eine Rechnung
des ersten Scanpunkts. Schlägt sie fehl, erscheint die Meldung von \FG; der vermutlich betroffene Reiter (bzw.\ die
Schicht) zeigt sie ebenfalls an.

\begin{tipp}[Vom Editor zum Skript]
Laden Sie die Konfiguration unter \gui{JSON} herunter. Dieselbe Datei rechnet \cmd{fgsolve datei.json} oder
\texttt{fg.solve(fg.load\_config("datei.json"))}.
\end{tipp}

%-------------------------------------------------------------------------------------------
\section{Einzelrechnung}
\begin{enumerate}
  \item \gui{Quelle}: \gui{Beispiel} oder \gui{Datei hochladen}.
  \item Die Kurzübersicht zeigt Dimension, Löser, Schichten und Materialien; unter \gui{JSON} steht die ganze
        Konfiguration.
  \item \gui{Rechnen} rechnet den ersten Scanpunkt (wie \texttt{fg.solve}).
  \item Ergebnis: $\sum R$, $\sum T$, Absorption, Energiebilanz, Tabelle aller propagierenden Ordnungen mit s/p-Anteil,
        Balkendiagramm, Warnungen und Diagnose (Laufzeiten, Genauigkeit, Konditionszahlen). \gui{Tabelle als CSV}
        lädt die Tabelle herunter.
\end{enumerate}

%-------------------------------------------------------------------------------------------
\section{Scan}
Rechnet die Konfiguration aus dem Editor vollständig: alle Werte von $\lambda$, $\varphi$, $\psi$, $\delta$ und den
\key{sweep}-Block (im Reiter \gui{JSON} des Editors anlegen). \gui{Scan rechnen} startet den Job; jeder Lauf bleibt in
der Sitzung und lässt sich auswählen.
\begin{itemize}
  \item \gui{x-Achse}: eine Größe, die im Scan mehrere Werte hat. Die übrigen veränderlichen Größen halten Sie fest.
  \item \gui{Ordnungen} $R(n,m)$, $T(n,m)$ und \gui{Summen} $\sum R$, $\sum T$; \gui{Größe}: gesamt, s (TE) oder p (TM).
        Eine Ordnung, die bei einem Wert nicht propagiert, steht als 0.
  \item \gui{Vergleichen mit}: weitere Läufe derselben Sitzung im selben Diagramm (gleiche $x$-Achse), z.\,B.\ zwei
        Werte von $N$ für eine Konvergenzprüfung.
  \item \gui{efficiencies.csv}: alle Zeilen wie bei \cmd{fgsolve}.
\end{itemize}

%-------------------------------------------------------------------------------------------
\section{Nahfeld}
Feldschnitte für den ersten Scanpunkt der Konfiguration aus dem Editor.
\begin{enumerate}
  \item \gui{Nahfeld-Einstellungen}: Fehlt der Block \key{nearfield}, legt \gui{Block mit Beispielwerten anlegen} ihn an.
        Dann Lage der Schnitte ($y_0$, $z_0$), $y$-Bereich, Punktzahlen, Perioden, Größen und Rekonstruktion einstellen.
  \item \gui{Nahfeld rechnen} startet den Job (\key{enabled} wird dafür auf \val{true} gesetzt).
  \item \gui{Schnitt}, \gui{Größe} ($|E|^2$, $|H|^2$, $S_y$, $E_x\ldots H_z$) und \gui{Anteil} (Betrag, Re, Im, Phase)
        wählen. Karten zeigen die Materialgrenzen als weiße Linien; 1D-Linien erscheinen als Kurve.
  \item \gui{Alle Schnitte (.zip)} enthält dieselben Dateien wie \cmd{fgsolve} (Abschnitt~\ref{sec:nf-dateien}).
\end{enumerate}

%-------------------------------------------------------------------------------------------
\section{CAD}
Für Schichten vom Typ \val{cad} (Kapitel~\ref{ch:cad}); setzt einen Server mit CAD-Slicer voraus (Seite \gui{System}:
„CAD-Slicer“).
\begin{enumerate}
  \item \gui{CAD-Schicht hinzufügen}, falls noch keine existiert.
  \item STEP-Datei hochladen (\file{.step}/\file{.stp}) oder ein Beispiel verwenden. Uploads liegen auf dem Server
        unter \file{uploads/} des Datenverzeichnisses; aus anderen Verzeichnissen liest die GUI nicht.
  \item \gui{Körper $\to$ Material}: Die Körpernamen werden aus der Datei gelesen; ordnen Sie jedem ein Material zu.
        Weitere Einstellungen: Einheit, Achsen, $y$-Bereich, Zellursprung, Hintergrund, Überlappung, Bandzerlegung,
        Normalenfeld.
  \item \gui{Prüfen} startet die Geometrieprüfung und zeigt Befunde (Code, Meldung, Ort in \textmu m), die Bänder mit
        Flächenanteilen und eine \textbf{Schnittvorschau} je Band oder bei einer eigenen Höhe.
\end{enumerate}
Danach rechnen Einzelrechnung, Scan oder Nahfeld die Konfiguration wie jede andere.

%-------------------------------------------------------------------------------------------
\section{Optimierung}\label{sec:gui-opt}
Maximiert oder minimiert $L=\sum w_ie_i$ mit L-BFGS-B und exakten Gradienten (Kapitel~\ref{ch:optimierung}) am ersten
Scanpunkt.
\begin{enumerate}
  \item \gui{Zielfunktion}: Tabelle aus Seite (R/T), Ordnung $n$, $m$ und Gewicht; mehrere Zeilen erlaubt.
        \gui{Ziel}: Maximieren oder Minimieren.
  \item \gui{Parameter}: JSON-Pointer wählen (\val{/layers/0/thickness\_um}, \val{/layers/0/breakpoints\_um/1},
        \val{/materials/al/n/1}, \val{/incidence/lambda\_um} \ldots), je Parameter Start, Unter- und Obergrenze. Die
        Vorgabe $\pm50$\,\% ist nur ein Platzhalter; der Start muss innerhalb der Grenzen liegen.
  \item \gui{Optimieren} startet den Job; der Verlauf erscheint live. \gui{Stoppen} beendet nach der laufenden
        Auswertung und behält das beste Ergebnis.
  \item Ergebnis: Zielwert am Start und am Optimum, Parameter, Verlauf (auch als CSV). \gui{Optimum in den Editor
        übernehmen} schreibt die Werte in die Konfiguration.
\end{enumerate}
\gui{Littrow-Vorlage laden} lädt das Beispiel \file{optimize\_littrow.py} (Aluminiumgitter, Periode 0,5\um, TM: Tiefe
und Stegbreite für die Ordnung $-1$).

%-------------------------------------------------------------------------------------------
\section{Rechnungen}
Alle Rechnungen (Einzelrechnung, Scan, Nahfeld) werden auf dem Server gespeichert und
bleiben nach einem Neustart in der Liste. Die Tabelle zeigt Art, Status, Platz in der Warteschlange und Laufzeit.
Für eine einzelne Rechnung:
\begin{itemize}
  \item \gui{Abbrechen} -- nur für wartende Rechnungen; eine laufende lässt sich nicht unterbrechen.
  \item \gui{Löschen} -- entfernt eine beendete Rechnung samt Dateien.
  \item \gui{Rechnung herunterladen (.zip)} -- Konfiguration, Ergebnistabellen (CSV/JSON, beim Nahfeld \file{.npy})
        und Rohdateien; \gui{Modell herunterladen (.json)} -- nur die Konfiguration.
  \item \gui{Im Editor öffnen} übernimmt die gespeicherte Konfiguration; \gui{Im Scan / Nahfeld öffnen} zeigt das
        Ergebnis dort.
\end{itemize}
Was bei einem Neustart des Servers noch wartete oder lief, steht als Fehler „Server wurde neu gestartet“ in der Liste;
bitte neu einreichen. Der Betreiber kann alte Rechnungen automatisch löschen lassen.

%-------------------------------------------------------------------------------------------
\section{System}
Version und Kern von \FG, ob CUDA nutzbar und der CAD-Slicer eingebaut ist, Backend-Voreinstellung,
GPU-Zustand (Speicher, Auslastung, Temperatur) und die Einstellungen der GUI. \gui{Neu prüfen} liest die Angaben
erneut. Der Abschnitt „Betrieb“ zeigt die Gesundheitsprüfung, den Speicherbedarf und \gui{Jetzt aufräumen}.

%-------------------------------------------------------------------------------------------
\section{Für Betreiber: Installation und Betrieb}\label{sec:gui-betrieb}
Die vollständige Anleitung steht in \file{GUI-FourierGrating/docs/deployment\_rtx6000.md}. Kurzfassung:
\begin{lstlisting}[style=shell]
sudo useradd --system --create-home --home-dir /opt/fg-gui --shell /sbin/nologin fggui
sudo -u fggui git clone https://github.com/PhotonPower/FourierGrating.git     /opt/fg-gui/FourierGrating
sudo -u fggui git clone https://github.com/PhotonPower/GUI-FourierGrating.git /opt/fg-gui/GUI-FourierGrating
sudo -u fggui /opt/fg-gui/GUI-FourierGrating/scripts/install_server.sh --prefix /opt/fg-gui
sudo cp deploy/fg-gui.env.example /etc/fg-gui/fg-gui.env      # anpassen
sudo cp deploy/fg-gui.service /etc/systemd/system/
sudo systemctl enable --now fg-gui
\end{lstlisting}
Ohne root-Rechte: \file{scripts/install\_user\_service.sh} richtet einen systemd-Benutzerdienst ein.

\begin{center}
\small
\begin{tabularx}{\linewidth}{lYl}
\toprule
Variable & Bedeutung & Standard \\
\midrule
\texttt{FG\_GUI\_DATA\_DIR} & Laufzeitdaten (Rechnungen, Uploads, Sitzungen) & \file{data} \\
\texttt{FG\_GUI\_EXAMPLES\_DIR} & Verzeichnis mit Beispielkonfigurationen & -- \\
\texttt{FG\_GUI\_WORKERS} & gleichzeitige Rechnungen & 1 \\
\texttt{FG\_GUI\_JOB\_HISTORY} & Rechnungen im Speicher & 50 \\
\texttt{FG\_GUI\_USER\_HEADER} & Header mit dem vom Proxy angemeldeten Benutzer (z.\,B.\ \texttt{X-Remote-User}) & -- \\
\texttt{FG\_GUI\_RETENTION\_DAYS} & beendete Rechnungen nach so vielen Tagen löschen & 0 (aus) \\
\texttt{FG\_GUI\_MAX\_RUNS} & höchstens so viele Rechnungen behalten & 0 (unbegrenzt) \\
\texttt{FG\_GUI\_LOG\_LEVEL}, \texttt{FG\_GUI\_LOG\_FILE} & Protokollstufe, zusätzliche Logdatei & \texttt{INFO} \\
\texttt{FG\_NUM\_THREADS}, \texttt{FG\_CORE}, \texttt{FG\_PRECISION}, \texttt{FG\_GPU\_STREAMS} & an \FG{} durchgereicht & \\
\bottomrule
\end{tabularx}
\end{center}
\paragraph{Zugriff.} Entweder über einen SSH-Tunnel
(\texttt{ssh -N -L 8505:server:8505 benutzer@server}, dann \texttt{http://localhost:8505}) oder über nginx mit TLS und
Anmeldung (Vorlage \file{deploy/nginx-fg-gui.conf}). Den Streamlit-Port \textbf{nicht} nach außen öffnen: Streamlit
hat keine eigene Anmeldung.
\paragraph{Überwachung.} \texttt{fg-gui-health --url \ldots} liefert Exit-Code 0, wenn die GUI gesund ist
(Nagios/Icinga/cron); das Protokoll steht im Journal (\texttt{journalctl -u fg-gui -f}).

\begin{achtung}[Was die GUI (noch) nicht kann]
\begin{itemize}[nosep]
  \item Dispersive Materialien: \FG{} kann sie seit M22, die GUI hat dafür noch kein Formular (geplant als G8).
  \item Eine laufende Rechnung abbrechen.
  \item Isolierte Strukturen (PML), Zylinder und ASR haben keine eigenen Formulare; sie werden im Reiter \gui{JSON}
        bearbeitet und funktionieren dann normal.
\end{itemize}
\end{achtung}
